% Generated from tex/chapters/ch04/05_ソフトウェア構築の基礎.tex for the SWEBOK structure tree
\subsubsection{変化を予期し受け入れる}

多くのソフトウェアは時間とともに変わる。要求が変わることもあれば、利用者数、実行環境、ライブラリ、法規制、ビジネス方針が変わることもある。構築で変化を予期するとは、将来の変更が入ったときに、局所的な修正で済むように作ることである。

変化を予期するためには、拡張しやすい構造、明確なインタフェース、設定による切り替え、再利用可能な部品、テストによる安全網が必要である。一方で、まだ必要でない過剰な柔軟性を作り込むと、逆に複雑さが増える。したがって、予期すべき変化を見極め、過剰設計を避けることが重要である。

現代の開発では、変化を避けるのではなく受け入れる実践が広く使われる。アジャイル開発、DevOps、継続的デリバリ、継続的デプロイは、短い周期で変更を入れ、テストし、提供し、運用から学ぶための仕組みである。
